\section{Technical Expertise}

% For me, every paper is a new opportunity to code something up. Hence, over the years, I have developed new techniques that are waiting to be employed at larger scales. Below, I summarize the technical expertise and tools I have developed over the years:

    \resumeSubHeadingListStart
      \resumeProjectHeading
          {\textbf{Observational} $|$ \emph{(Sub)mm-wave interferometry, bolometry, CMB-surveys}}{}
          \resumeItemListStart
            \resumeItem{I am acquainted with using multiwavelength observations, particularly \textit{HST} observations, \textit{Chandra} imaging, FIR interferometry (ALMA) as well as bolometry (ACT) data, and radio surveys.}
            \resumeItem{I am experienced with spectroscopy as well as brought-band continuum imaging, handling time-ordered data, and map-making/component separation.}
            
          \resumeItemListEnd
          
      \resumeProjectHeading
          {\textbf{Fourier domain} $|$ \emph{Visibility modeling, Nested sampling, Matched filters, Multiprocessing}}{}
          \resumeItemListStart
            \resumeItem{I worked on the development of a generic $uv$-plane based Bayesian forward modeling routine.}
            % \resumeItem{I wrote a routine that models frequency maps directly (instead of the more commonly used Compton-Y maps).}
            \resumeItem{I specialized in joint-likelihood modeling, as well as writing modules for data combinations of single-dish data and interferometric observations.}
            \resumeItem{I wrote a new interferometric deconvolution algorithm optimized for imaging large spatial scales.}
            \resumeItem{I implemented a matched filter in the \textit{uv}-plane to detect weak lines in interferometric data.}
            \resumeItem{I made a routine that by Jackknifing the visibilities creates ``empty'' noise cubes for ALMA observations while accounting for time-sensitive observing conditions.}
          \resumeItemListEnd

    \resumeProjectHeading
          {\textbf{Simulational} $|$ \emph{Hydrodynamical simulations, phenomenological models, atmospheric modeling}}{}
          \resumeItemListStart
            \resumeItem{I wrote a semi-analytical phenomenological model that creates mock X-ray observations of galaxy clusters from cosmological hydrodynamical simulations such as the BAHAMAS suite.}
            \resumeItem{I wrote a tool to create synthetic bolometric FIR observations using realistic atmospheric models used to optimize scanning strategies and bridge the gap between simulations and observations.}
          \resumeItemListEnd

    % \resumeProjectHeading
    %       {\textbf{Tensorflow} $|$ \emph{CNNs, RNNs, GANs}}{}
    %       \resumeItemListStart
    %         \resumeItem{I followed classes, getting into the details of neural networks and other machine learning techniques like PCA.}
    %       \resumeItemListEnd
    \resumeSubHeadingListEnd
